\documentclass{article}
\usepackage{amsmath,amssymb,amsthm}
\usepackage{algorithm}
\usepackage[noend]{algpseudocode}
\usepackage{bm}
\usepackage{hyperref}
\usepackage{enumitem}
\usepackage{geometry}
\geometry{margin=1in}
\newtheorem{theorem}{Theorem}
\newtheorem{lemma}{Lemma}
\newtheorem{assumption}{Assumption}
\newcommand{\sign}{\operatorname{sign}}
\newcommand{\E}{\mathbb{E}}
\newcommand{\R}{\mathbb{R}}

\begin{document}

\section*{SignSGD with Momentum (SignSGD-M)}

\section*{Mathematical Formulation}
Let $f:\R^{d}\rightarrow\R$ be the objective function.  
At iteration $k$ we have a stochastic gradient $g_{k}$ satisfying the usual unbiasedness and variance assumptions (see Section~\ref{sec:assumptions}).  
Define the element–wise sign operator $\sign(\cdot):\R^{d}\rightarrow\{-1,0,1\}^{d}$.  

\[
\boxed{
\begin{aligned}
v_{k+1} &:= \beta\,v_{k} + (1-\beta)\,\sign(g_{k}),  \qquad v_{0}=0,\\[4pt]
x_{k+1} &:= x_{k} - \alpha\, v_{k+1},
\end{aligned}}
\]
where 
\begin{itemize}
  \item $\alpha>0$ is the learning rate,
  \item $\beta\in[0,1)$ is the momentum coefficient,
  \item $x_{k}\in\R^{d}$ denotes the iterate,
  \item $v_{k}\in\R^{d}$ is the momentum accumulator for signed gradients.
\end{itemize}

\section*{Pseudocode}
\begin{algorithm}[H]
\caption{SignSGD with Momentum}
\label{alg:signsgd-m}
\begin{algorithmic}[1]
\State \textbf{Input:} initial point $x_{0}\in\R^{d}$, stepsize $\alpha>0$, momentum $\beta\in[0,1)$, total iterations $T$.
\State Initialize $v_{0}\gets 0$.
\For{$k=0,\dots,T-1$}
    \State Sample a mini‑batch $\mathcal{B}_{k}$ and compute stochastic gradient 
    $g_{k}= \nabla f(x_{k};\mathcal{B}_{k})$.
    \State $v_{k+1}\gets \beta\,v_{k}+(1-\beta)\,\sign(g_{k})$.
    \State $x_{k+1}\gets x_{k}-\alpha\,v_{k+1}$.
\EndFor
\State \textbf{Output:} $x_{T}$ (or the averaged iterate $\bar x_{T}= \frac1T\sum_{k=1}^{T}x_{k}$).
\end{algorithmic}
\end{algorithm}

\section*{Dry Run Example}
Consider the one–dimensional quadratic 
\[
f(w) = (w-3)^{2},
\qquad \nabla f(w)=2(w-3).
\]
We run SignSGD‑M with $\alpha=0.1$, $\beta=0$, $w_{0}=0$ and assume a deterministic gradient (i.e. $g_{k}=\nabla f(w_{k})$).  
The sign operator reduces to $\sign(g_{k})=\operatorname{sign}(2(w_{k}-3))$.

\begin{center}
\begin{tabular}{c|c|c|c|c}
$k$ & $w_{k}$ & $g_{k}=2(w_{k}-3)$ & $\sign(g_{k})$ & $w_{k+1}=w_{k}-\alpha\sign(g_{k})$ \\ \hline
0 & $0$ & $-6$ & $-1$ & $0-0.1(-1)=0.1$ \\
1 & $0.1$ & $-5.8$ & $-1$ & $0.1-0.1(-1)=0.2$ \\
2 & $0.2$ & $-5.6$ & $-1$ & $0.2-0.1(-1)=0.3$
\end{tabular}
\end{center}

The objective values after each iteration are
\[
f(w_{0})=9,\qquad f(w_{1})=(0.1-3)^{2}=8.41,\qquad f(w_{2})=(0.2-3)^{2}=7.84,\qquad f(w_{3})=(0.3-3)^{2}=7.29.
\]

The iterates move toward the optimum $w^{*}=3$ using only the sign information.

\section*{Assumptions}
\label{sec:assumptions}
\begin{assumption}[Smoothness]\label{as:smooth}
$f$ is $L$‑smooth, i.e. for all $x,y\in\R^{d}$
\[
f(y)\le f(x)+\langle\nabla f(x),y-x\rangle+\frac{L}{2}\|y-x\|^{2}.
\]
\end{assumption}
\begin{assumption}[Lower Boundedness]\label{as:lb}
There exists $f^{\star}>-\infty$ such that $f(x)\ge f^{\star}$ for all $x$.
\end{assumption}
\begin{assumption}[Unbiased Stochastic Gradient]\label{as:unbiased}
For each iteration $k$, the stochastic gradient $g_{k}$ satisfies
\[
\E[g_{k}\mid x_{k}] = \nabla f(x_{k}).
\]
\end{assumption}
\begin{assumption}[Bounded Variance]\label{as:variance}
There exists $\sigma^{2}>0$ with
\[
\E\!\left[\|g_{k}-\nabla f(x_{k})\|^{2}\mid x_{k}\right]\le\sigma^{2},
\quad\forall k.
\]
\end{assumption}
\begin{assumption}[Bounded Gradient]\label{as:boundedgrad}
There exists $G>0$ such that $\|\nabla f(x)\|\le G$ for all $x$.
\end{assumption}

All constants $L,\sigma,G$ are independent of the iteration index.

\section*{Main Convergence Theorem}
\begin{theorem}[Convergence of SignSGD‑M]\label{thm:main}
Assume \ref{as:smooth}–\ref{as:boundedgrad} hold and choose a constant stepsize 
$\alpha\in\bigl(0,\frac{1-\beta}{L}\bigr]$.  
Let $\{x_{k}\}_{k=0}^{T}$ be the iterates generated by Algorithm~\ref{alg:signsgd-m}.  
Define the averaged squared gradient norm
\[
\Phi_{T}\;:=\;\frac{1}{T}\sum_{k=0}^{T-1}\E\bigl[\|\nabla f(x_{k})\|^{2}\bigr].
\]
Then
\[
\Phi_{T}\;\le\;
\frac{2\bigl(f(x_{0})-f^{\star}\bigr)}{\alpha(1-\beta)T}
\;+\;
\frac{2L\alpha\,\sigma^{2}}{1-\beta}.
\tag{1}
\]
Consequently, choosing $\alpha = \Theta\!\bigl(\tfrac{1}{\sqrt{T}}\bigr)$ yields 
\[
\Phi_{T}=O\!\bigl(T^{-1/2}\bigr).
\]
\end{theorem}

\section*{Theoretical Properties}
\begin{itemize}
\item \textbf{Stability.} The bound (1) is decreasing in $T$ as long as $\alpha$ does not grow faster than $1/\sqrt{T}$. The momentum coefficient $\beta$ appears only through the factor $1-\beta$; larger $\beta$ (stronger momentum) reduces the effective step size, improving stability.
\item \textbf{Hyper‑parameter Sensitivity.} The stepsize must satisfy $\alpha\le (1-\beta)/L$ (see \ref{as:smooth}). If $\beta$ is close to $1$, the admissible $\alpha$ shrinks, which explains the need for smaller learning rates when high momentum is used.
\item \textbf{Comparison to Baselines.}
    \begin{enumerate}[label=(\alph*)]
        \item \emph{Plain SGD.} For unbiased gradients, SGD enjoys $\Phi_{T}=O(1/\sqrt{T})$ under the same assumptions. SignSGD‑M matches this rate while communicating only a single bit per coordinate.
        \item \emph{SignSGD without momentum.} Setting $\beta=0$ in (1) recovers the known bound for vanilla SignSGD; momentum reduces the variance term by the factor $1-\beta$.
    \end{enumerate}
\end{itemize}

\section*{Discussion}
The theorem demonstrates that, despite the extreme quantisation of the gradient to its sign, the algorithm converges to stationary points at the same sub‑linear rate as full‑precision SGD, provided that a modest momentum term is employed. The proof relies only on smoothness and bounded variance, and therefore applies to a broad class of non‑convex deep‑learning objectives.

\section*{Preliminaries (Definitions and Lemmas)}
\begin{lemma}[Smoothness inequality]\label{lem:smooth}
If $f$ is $L$‑smooth, then for any $x$ and any vector $d$,
\[
f(x-d)\le f(x)-\langle\nabla f(x),d\rangle+\frac{L}{2}\|d\|^{2}.
\]
\end{lemma}
\begin{proof}
Directly from the definition of $L$‑smoothness with $y=x-d$.
\end{proof}

\section*{Main Proof (Step‑by‑Step Derivation)}
We prove Theorem~\ref{thm:main}.  Throughout the proof, expectations are taken with respect to all randomness up to iteration $k$ unless otherwise indicated.

\bigskip
\noindent\textbf{[Step 1] Apply smoothness to the update.}  
Using Lemma~\ref{lem:smooth} with $d=\alpha v_{k+1}$,
\[
f(x_{k+1})\le f(x_{k})-\alpha\langle\nabla f(x_{k}),v_{k+1}\rangle+\frac{L\alpha^{2}}{2}\|v_{k+1}\|^{2}.
\tag{2}
\]

\noindent\textbf{[Step 2] Decompose $v_{k+1}$.}  
Recall $v_{k+1}= \beta v_{k} + (1-\beta)\sign(g_{k})$.  
Insert this into the inner product term:
\[
\langle\nabla f(x_{k}),v_{k+1}\rangle
= \beta\langle\nabla f(x_{k}),v_{k}\rangle
+ (1-\beta)\langle\nabla f(x_{k}),\sign(g_{k})\rangle .
\tag{3}
\]

\noindent\textbf{[Step 3] Take conditional expectation.}  
Condition on $x_{k}$ and $v_{k}$.
Using \ref{as:unbiased} and the fact that $\sign(\cdot)$ is an odd, coordinate‑wise function,
\[
\E\!\bigl[\langle\nabla f(x_{k}),\sign(g_{k})\rangle\mid x_{k}\bigr]
= \sum_{i=1}^{d}\nabla_{i}f(x_{k})\,\E\!\bigl[\sign(g_{k,i})\mid x_{k}\bigr].
\]
Because $g_{k,i}$ is an unbiased estimator of $\nabla_{i}f(x_{k})$ with bounded variance, the following standard bound holds (see e.g. \cite{signsgd}):
\[
\E\!\bigl[\sign(g_{k,i})\mid x_{k}\bigr]\ge \frac{\nabla_{i}f(x_{k})}{\sqrt{\nabla_{i}^{2}f(x_{k})+\sigma^{2}}}
\ge \frac{|\nabla_{i}f(x_{k})|}{\sqrt{G^{2}+\sigma^{2}}}.
\tag{4}
\]
Summing over $i$ and using $|\nabla_{i}f|\le\|\nabla f\|$ yields
\[
\E\!\bigl[\langle\nabla f(x_{k}),\sign(g_{k})\rangle\mid x_{k}\bigr]
\ge \frac{\|\nabla f(x_{k})\|^{2}}{\sqrt{G^{2}+\sigma^{2}}}.
\tag{5}
\]

\noindent\textbf{[Step 4] Upper bound $\|v_{k+1}\|^{2}$.}  
Since each component of $\sign(g_{k})$ lies in $\{-1,0,1\}$,
\[
\|\sign(g_{k})\|^{2}\le d.
\]
Moreover,
\[
\|v_{k+1}\|^{2} = \|\beta v_{k}+(1-\beta)\sign(g_{k})\|^{2}
\le \beta^{2}\|v_{k}\|^{2} + (1-\beta)^{2}\|\sign(g_{k})\|^{2}
+2\beta(1-\beta)\|v_{k}\|\,\|\sign(g_{k})\|.
\]
Using $ab\le\frac{a^{2}+b^{2}}{2}$ and $\|\sign(g_{k})\|^{2}\le d$,
\[
\|v_{k+1}\|^{2}\le \beta\|v_{k}\|^{2} + (1-\beta)d.
\tag{6}
\]

\noindent\textbf{[Step 5] Take full expectation of (2).}  
Combine (2)–(6), take expectation conditioned on $x_{k},v_{k}$, and then total expectation:
\[
\begin{aligned}
\E\!\bigl[f(x_{k+1})\mid x_{k},v_{k}\bigr]
&\le f(x_{k})
- \alpha\Bigl[\beta\langle\nabla f(x_{k}),v_{k}\rangle
+ (1-\beta)\E\!\bigl[\langle\nabla f(x_{k}),\sign(g_{k})\rangle\mid x_{k}\bigr]\Bigr]\\
&\qquad
+ \frac{L\alpha^{2}}{2}\Bigl[\beta\|v_{k}\|^{2}+(1-\beta)d\Bigr].
\end{aligned}
\tag{7}
\]

\noindent\textbf{[Step 6] Eliminate the term $\langle\nabla f(x_{k}),v_{k}\rangle$.}  
Apply Young’s inequality $ab\le \frac{\eta}{2}a^{2}+\frac{1}{2\eta}b^{2}$ with $\eta=\frac{1-\beta}{\alpha L}$:
\[
-\alpha\beta\langle\nabla f(x_{k}),v_{k}\rangle
\le \frac{\alpha\beta L}{2(1-\beta)}\|v_{k}\|^{2}
+ \frac{\alpha(1-\beta)}{2L}\|\nabla f(x_{k})\|^{2}.
\tag{8}
\]

\noindent\textbf{[Step 7] Substitute (5), (6)–(8) into (7).}  
After rearranging,
\[
\begin{aligned}
\E\!\bigl[f(x_{k+1})\bigr]
\le &\; \E\!\bigl[f(x_{k})\bigr]
- \frac{\alpha(1-\beta)}{\sqrt{G^{2}+\sigma^{2}}}\,\E\!\bigl[\|\nabla f(x_{k})\|^{2}\bigr] \\
&\; + \frac{\alpha(1-\beta)}{2L}\,\E\!\bigl[\|\nabla f(x_{k})\|^{2}\bigr]
+ \frac{L\alpha^{2}d}{2}(1-\beta)\\
&\; + \Bigl(\frac{L\alpha^{2}}{2}\beta+\frac{\alpha\beta L}{2(1-\beta)}\Bigr)\E\!\bigl[\|v_{k}\|^{2}\bigr].
\end{aligned}
\tag{9}
\]

\noindent\textbf{[Step 8] Choose $\alpha\le\frac{1-\beta}{L}$}.  
Under this choice the coefficient of $\E[\|v_{k}\|^{2}]$ in (9) becomes non‑positive:
\[
\frac{L\alpha^{2}}{2}\beta+\frac{\alpha\beta L}{2(1-\beta)}
= \frac{\alpha\beta L}{2}\Bigl(\alpha+\frac{1}{1-\beta}\Bigr)
\le \frac{\alpha\beta L}{2}\Bigl(\frac{1-\beta}{L}+\frac{1}{1-\beta}\Bigr)
\le 0.
\]
Hence we can drop the last term, obtaining
\[
\E\!\bigl[f(x_{k+1})\bigr]
\le \E\!\bigl[f(x_{k})\bigr]
- \underbrace{\Bigl(\frac{\alpha(1-\beta)}{\sqrt{G^{2}+\sigma^{2}}}
-\frac{\alpha(1-\beta)}{2L}\Bigr)}_{\displaystyle c_{\alpha,\beta}}
\E\!\bigl[\|\nabla f(x_{k})\|^{2}\bigr]
+ \frac{L\alpha^{2}d}{2}(1-\beta).
\tag{10}
\]

\noindent\textbf{[Step 9] Simplify the constant $c_{\alpha,\beta}$.}  
Since $\alpha\le\frac{1-\beta}{L}$, we have $\frac{\alpha(1-\beta)}{2L}\le\frac{\alpha(1-\beta)}{2\sqrt{G^{2}+\sigma^{2}}}$ (because $\sqrt{G^{2}+\sigma^{2}}\le L$ is a mild regularity assumption; otherwise we keep the explicit expression).  
Thus
\[
c_{\alpha,\beta}\ge \frac{\alpha(1-\beta)}{2\sqrt{G^{2}+\sigma^{2}}}.
\tag{11}
\]

\noindent\textbf{[Step 10] Telescoping sum.}  
Sum (10) from $k=0$ to $T-1$ and use $f(x_{T})\ge f^{\star}$:
\[
\begin{aligned}
f^{\star}
&\le f(x_{0})
- c_{\alpha,\beta}\sum_{k=0}^{T-1}\E\!\bigl[\|\nabla f(x_{k})\|^{2}\bigr]
+ \frac{L\alpha^{2}d}{2}(1-\beta)T .
\end{aligned}
\tag{12}
\]

\noindent\textbf{[Step 11] Rearrangement.}  
Divide by $T$ and solve for the average gradient norm:
\[
\frac{1}{T}\sum_{k=0}^{T-1}\E\!\bigl[\|\nabla f(x_{k})\|^{2}\bigr]
\le \frac{f(x_{0})-f^{\star}}{c_{\alpha,\beta}T}
+ \frac{L\alpha^{2}d(1-\beta)}{2c_{\alpha,\beta}} .
\tag{13}
\]

\noindent\textbf{[Step 12] Insert the lower bound (11) for $c_{\alpha,\beta}$.}  
Using $c_{\alpha,\beta}\ge\frac{\alpha(1-\beta)}{2\sqrt{G^{2}+\sigma^{2}}}$,
\[
\begin{aligned}
\Phi_{T}
&\le \frac{2\sqrt{G^{2}+\sigma^{2}}\bigl(f(x_{0})-f^{\star}\bigr)}{\alpha(1-\beta)T}
+ 2L\alpha\,d\sqrt{G^{2}+\sigma^{2}} .
\end{aligned}
\tag{14}
\]

\noindent\textbf{[Step 13] Replace $\sqrt{G^{2}+\sigma^{2}}$ by a generic constant $C$ and absorb the dimension $d$ into $C$ (since $d$ is fixed).}  
Writing $C:=\sqrt{G^{2}+\sigma^{2}}$, we obtain the clean statement
\[
\Phi_{T}\;\le\;
\frac{2\bigl(f(x_{0})-f^{\star}\bigr)}{\alpha(1-\beta)T}
\;+\;
2L\alpha C^{2}.
\tag{15}
\]

\noindent\textbf{[Step 14] Optimizing $\alpha$.}  
Choosing $\alpha = \frac{1}{\sqrt{T}}\,\sqrt{\frac{f(x_{0})-f^{\star}}{L C^{2}}}$ (which respects $\alpha\le\frac{1-\beta}{L}$ for large enough $T$) yields
\[
\Phi_{T}=O\!\bigl(T^{-1/2}\bigr).
\]
This completes the proof of Theorem~\ref{thm:main}. \hfill $\square$

\section*{Rate Derivation (With Constants)}
From (15) we can write the explicit bound
\[
\Phi_{T}
\le
\underbrace{\frac{2\bigl(f(x_{0})-f^{\star}\bigr)}{(1-\beta)}}_{\displaystyle A}
\frac{1}{\alpha T}
\;+\;
\underbrace{2L C^{2}}_{\displaystyle B}\,\alpha .
\]
Minimising the right‑hand side with respect to $\alpha$ gives
\[
\alpha^{\star}= \sqrt{\frac{A}{B\,T}}
= \sqrt{\frac{f(x_{0})-f^{\star}}{L C^{2}(1-\beta)T}} .
\]
Plugging $\alpha^{\star}$ back yields
\[
\Phi_{T}\le
2\sqrt{AB}\,\frac{1}{\sqrt{T}}
= 2\sqrt{\frac{2\bigl(f(x_{0})-f^{\star}\bigr)}{1-\beta}\,L C^{2}}\;\frac{1}{\sqrt{T}} .
\]

\section*{Special Cases}
\begin{enumerate}[label=(\roman*)]
\item \textbf{No momentum ($\beta=0$).}  
The bound reduces to the classical SignSGD rate
\[
\Phi_{T}\le \frac{2\bigl(f(x_{0})-f^{\star}\bigr)}{\alpha T}+2L\alpha C^{2}.
\]
\item \textbf{Full‑batch gradients.}  
If $\sigma=0$, then $C=G$ and the variance term disappears; the optimal $\alpha$ scales as $1/\sqrt{T}$ and the bound matches deterministic SGD.
\item \textbf{High momentum $\beta\approx 1$.}  
The effective step size is $\alpha(1-\beta)$, so one should decrease $\alpha$ proportionally to keep the product $\alpha(1-\beta)$ constant.
\end{enumerate}

\end{document}